% trainable-readout-numerics.tex -- gerado a mao a partir de readout_scale.log / readout_scale.jsonl
% (experiments/affine-forcing), run de 2026-07-27, 10.3 min, 2 workers nice -19.
% [RESULT 2026-07-27]

\paragraph{The readout freezes from the initialization --- a separate fact from
\eqref{eq:v0const}.}
Over the accepted instances,
\[
\max\bigl|v_0(\thit)-v_0(0)\bigr|
=1.2\cdot10^{-6}\ \ (v_0(0)=0.5),
\qquad
3.6\cdot10^{-6}\ \ (v_0(0)=2.0).
\]
\textbf{This is not a test of \eqref{eq:v0const}}, and we will not present it as
one. \Cref{eq:v0const} is a statement about the \emph{episode}
$[\thit,\trel]$, where the unit is dead; the measurement above is over
$[0,\thit]$, where the unit is \emph{alive}. It says that in this drawing family
$v_0$ barely moves even before the death, because $A$ is initialized with a
small margin on the face ($z_A(0)\le0.05$) and its output --- to which
$\partial L/\partial v_0$ is proportional --- is small throughout. It is a
welcome empirical fact, because it is what makes the discriminating design work
at all: it lets us set $v_0(\thit)$ by setting $v_0(0)$. Testing
\eqref{eq:v0const} properly would require logging $v_0$ \emph{inside} the
episode against a $\beta$-sweep, which is not done (\cref{sec:open}\,(P2)).
Practically, the consequence here is that the two predictors now differ by a
clean factor $2$ and $\tfrac12$. The instrument is no longer null.

\paragraph{The corrected law wins, and by the predicted factor.}
Acceptance: $22/30$ at $v_0(0)=0.5$ and $26/30$ at $v_0(0)=2.0$ (rejections: $11$
$\beta$-runs with no release inside the horizon, $1$ death before $t_{\min}$; no
instance was rejected by any readout sign guard, so \hyplabel{H6$'$} was never
binding). \textbf{The rejection rate depends on the treatment} --- $8/30$ against
$4/30$ --- so the two arms are not draws from a common population and any
\emph{between-arm} comparison of error levels is confounded with selection. We
therefore do not make one. The comparison that decides item~(a) is
\emph{within} each instance, where the two predictors are computed from the same
trajectory, and selection cannot touch it. Registered as
\cref{sec:open}\,(P9). Writing
$\mu=\beta(t^\beta_{\rm release}-\texit)$ for the measured coefficient:

\begin{center}
\begin{tabular}{@{}lcccc@{}}
\toprule
& \multicolumn{2}{c}{$v_0(0)=0.5$ ($N=22$)} & \multicolumn{2}{c}{$v_0(0)=2.0$ ($N=26$)} \\
\cmidrule(lr){2-3}\cmidrule(lr){4-5}
& $\mu/\text{naive}$ & $\mu/\text{corr.}$ & $\mu/\text{naive}$ & $\mu/\text{corr.}$ \\
\midrule
predicted by \eqref{eq:release} & $2.000$ & $1.000$ & $0.500$ & $1.000$ \\
\textbf{measured, median}       & $\mathbf{2.049}$ & $\mathbf{1.024}$ & $\mathbf{0.496}$ & $\mathbf{0.991}$ \\
measured, range                 & $[1.63,2.95]$ & $[0.81,1.47]$ & $[0.48,0.69]$ & $[0.97,1.38]$ \\
\bottomrule
\end{tabular}
\end{center}

\noindent
The medians of the ratio to the \emph{naive} predictor are $2.049$ and $0.496$
against the arithmetic values $2$ and $\tfrac12$ that \eqref{eq:release}
predicts, with no free parameter and no fit. Relative errors:

\begin{center}
\begin{tabular}{@{}lccc@{}}
\toprule
 & median & p90 & max \\
\midrule
$|\mu-\text{corrected}|/|\text{corrected}|$, pooled $N=48$ & $\mathbf{2.76\%}$ & $17.3\%$ & $47.3\%$ \\
\ \ at $v_0(0)=0.5$ & $5.80\%$ & $16.7\%$ & $47.3\%$ \\
\ \ at $v_0(0)=2.0$ & $2.19\%$ & $15.3\%$ & $37.7\%$ \\
\midrule
$|\mu-\text{naive}|/|\text{naive}|$ (\emph{arithmetic floor}) & $51.4\%$ & --- & --- \\
\bottomrule
\end{tabular}
\end{center}

\noindent
\textbf{The last row is not an empirical error and must not be read as one.} If
\eqref{eq:release} is exactly right then, identically,
\[
\frac{|\mu-\text{naive}|}{|\text{naive}|}
\ \approx\ \Bigl|1-\frac1{v_0}\Bigr|
=\begin{cases}100\% & v_0=0.5\\ 50\% & v_0=2.0,\end{cases}
\]
which is what is observed ($104.9\%$ and $50.4\%$). It is an
\emph{arithmetic floor} built into the contrast by the choice of $v_0(0)$: it
measures how far apart we placed the two hypotheses, not how badly the naive one
fits. Reporting it as ``the naive law errs by $51\%$'' would be inflating a
design parameter into a result. \textbf{The informative statistics are the
paired count and the ratios $\mu/\text{predictor}$}, both of which are free of
this artefact.

\smallskip\noindent
\textbf{Paired verdict: the corrected predictor beats the naive one in $47$ of
$48$ instances} (one-sided sign test, $p=1.7\cdot10^{-13}$), with pooled median
ratio $\texttt{err\_corr}/\texttt{err\_naive}=\mathbf{0.046}$ --- a factor $22$.
(The value $0.043$ is the $v_0(0)=2.0$ subgroup, not the pool.) The single
exception is at $v_0(0)=2.0$ and is an instance whose corrected error is already
the sample maximum ($37.7\%$), i.e.\ an instance where the residual $\beta_0$
effect dominates the $v_0$ effect, not one where the naive law is right.

\smallskip\noindent
\emph{Independent replication.} Fresh seeds (\path{review18_rep.py},
\texttt{SEED0}$=77770101$, disjoint from the run above), $8$ instances per level,
$16/16$ accepted: $\mu/\text{naive}$ medians $1.950$ and $0.4995$,
$\mu/\text{corrected}$ medians $0.975$ and $0.999$, paired wins $15/16$, and
$\max|\Lambda_{\rm meas}-\Lambda|=7.1\cdot10^{-15}$ again. Nothing in the
verdict depends on the particular seed block.

\smallskip\noindent
\emph{The falsification condition fixed before the run was not met.} Item~(a) of
\cref{thm:d2} stands: the first-order coefficient carries the dead unit's
readout.

\paragraph{The lemma is again saturated, and Gr\"onwall is again not the driver.}
Measuring $\Lambda_{\rm meas}:=\max(0,-\lambda_{\min}(\nabla^2U))$ by exact
eigenvalues of the $12\times12$ Hessian along the hybrid trajectory, against
$\Lambda=\max_k\|c_k\|$: $\max_{48}|\Lambda_{\rm meas}-\Lambda|=7.1\cdot10^{-15}$,
pointwise relative $\max_{48}\bigl(|\Lambda_{\rm meas}-\Lambda|/\Lambda\bigr)
=1.59\cdot10^{-15}$. (The figure $3.0\cdot10^{-15}$ quoted elsewhere is the
different normalization $\max|\cdot|/\min\Lambda$, which is an upper bound on
the pointwise one; we report both and label them.) \emph{$\Lambda_{\rm meas}$ is
not the $L_B$ of \cite{paper}}, which is $\lambda_{\max}(G_B)$ --- see
\cref{rem:O6}; the whole point is that the two differ. Ranges: $\Lambda\in[2.36,7.07]$ (median $4.34$),
$\Tep\in[0.0145,0.351]$ (median $0.0626$), and the episode Gr\"onwall factor
\[
e^{\Lambda\Tep}\in[1.08,\ 2.87],\qquad\text{median }1.28 .
\]
The correlation of $\log\texttt{err\_corr}$ with $\log e^{\Lambda\Tep}$ is
$\mathbf{-0.45}$ --- \emph{negative}, as in \cite{scoping}, so the exponential
factor cannot be the causal channel for the residual error --- while the
correlation with $\log(1/R^{\hyb}(\texit))$ is $+0.29$, the same
loss-of-transversality driver that (O4) of \cite{paper} identifies for the
frozen readout. Both correlations are weaker than in \cite{scoping}
($-0.63$, $+0.60$), which is expected: pooling two $v_0$ levels mixes two
different drawing families. \emph{We do not claim a decomposition;} $N=48$ at a
single $\beta$ excludes Gr\"onwall as the dominant channel by the sign, and
nothing more.

\paragraph{Guard and balancedness.}
Hypothesis \hyplabel{H6$'$} as stated concerns $v_0$ only, and it is satisfied
\emph{trivially} here: $|v_0|$ stays within $4\cdot10^{-6}$ of $v_0(0)\in
\{0.5,2.0\}$, so $\varphi_v\in\{0.5,2.0\}$ effectively. The quantity
$\min_k\min_t|v_k|$ over the \emph{live} units, which the acceptance filter also
guards but which \hyplabel{H6$'$} no longer requires, has minimum
$5.0\cdot10^{-3}$ over the sample (median over instances $0.106$); no sign
crossing occurred inside a window in any accepted instance. Balancedness drift
$\max_k|B_k(t)-B_k(0)|$ over the episode: $\le6.0\cdot10^{-5}$ on the hybrid
($dt=10^{-5}$) and $\le9.5\cdot10^{-6}$ on the $\beta$-run
($dt=2\cdot10^{-6}$), a ratio $\approx6.3$ against a step ratio of $5$ --- the
drift is Euler error, not softplus defect, as \cref{sec:open}\,(P3) predicts.
